%% ED26-SCI-SaltelliAnnoni2010-001.
\chapter{Global Sensitivity and Robustness}
\label{sec:p15-global-sensitivity}

A model can be insensitive along the particular changes that were tried
and still change substantially when inputs vary together.
Sensitivity analysis must specify the output, the allowable inputs,
their joint distribution when probabilities are used, and the conclusion
being assessed. This entry reconstructs the argument and method choices
of Saltelli and Annoni \cite{SaltelliAnnoni2010Sensitivity}.
The variance decompositions and computation identities are derived below
from the conditional-expectation and product-measure foundations.
\index{sensitivity analysis!global}\index{one-at-a-time analysis}
\index{variance-based sensitivity}\index{robustness!model-based inference}

\section{The question and its domain}

\begin{hbdefinition}{def:p15-global-model}{Model, uncertainty law and target}
Let \(d\geq1\) be an integer, let \(\mathcal X\subseteq\mathbb R^d\)
be a Borel input domain with its trace sigma-algebra, and let
\(f:\mathcal X\to\mathbb R\) be a measurable
deterministic model output. When uncertainty is probabilistic, specify
a joint law \(P\) of \(X=(X_1,\ldots,X_d)\) supported on \(\mathcal X\)
and set \(Y=f(X)\). A baseline \(b\in\mathcal X\) is one input vector.

An uncertainty analysis concerns the distribution of \(Y\), such as
its variance, quantiles or probability of exceeding a threshold.
A sensitivity analysis attributes variation, or tests a conclusion,
relative to specified changes of inputs or assumptions.
A local derivative \(\partial_i f(b)\) concerns infinitesimal variation
at one point. A finite change \(f(b+he_i)-f(b)\) concerns a particular
coordinate segment. Neither object alone specifies the distribution
of the output over \(\mathcal X\).

A robustness claim additionally names its conclusion and alternatives:
for example, \(f(x)\leq c\) throughout a stated set, or a ranking that
persists over a stated family of input laws and model variants.
The range, law, dependence assumptions, output transformation and
threshold are part of the claim.
\end{hbdefinition}

Independent continuous input laws can be generated from independent
uniform coordinates by their quantile functions. Generalized quantiles
also accommodate atoms. Use interior probabilities where an endpoint
quantile is infinite. Under dependence, separate marginal transforms
do not define the joint law; a design must preserve the declared joint
distribution and feasibility constraints.
The measure used to weight an input region is distinct from its ordinary
geometric volume.

\section{What one-at-a-time exploration can establish}

\begin{hbdefinition}{def:p15-global-oat}{A fixed-baseline OAT design}
At baseline \(b\), a one-at-a-time design evaluates points for which
at most one coordinate differs from \(b\).
Its complete coordinate-line set is
\[
 A_b=\bigcup_{i=1}^d
 \{x\in\mathcal X:x_j=b_j\text{ for every }j\ne i\}.
\]
A finite OAT design uses a subset of this set.
An exact nonzero output difference establishes dependence along that
specific coordinate change. A zero difference does not establish
global irrelevance. With noisy or numerically inaccurate evaluations,
even the interpretation of a nonzero observed difference requires
an error analysis.
\end{hbdefinition}

\begin{hbtheorem}{thm:p15-global-oat-obstruction}
 {proved=annex:p15.global.oat-obstruction}{Coordinate lines do not identify interaction}
For \(d\geq2\), on any rectangle with nonzero widths about \(b\), the
functions
\[
 f_0(x)=0,\qquad f_C(x)=C(x_1-b_1)(x_2-b_2)
\]
agree on all of \(A_b\), for every real \(C\).
Nevertheless \(f_C\) has nonzero values outside \(A_b\)
when \(C\ne0\), and its magnitude at such points can be made arbitrarily large.
Thus even complete knowledge of every fixed-baseline coordinate curve
does not identify the global output of a general continuous model.
\end{hbtheorem}

\begin{hbtheorem}{thm:p15-global-ball-bound}
 {proved=annex:p15.global.ball-bound}{The centered-cube geometry}
In the normalized cube \([0,1]^d\) with baseline
\(b=(1/2,\ldots,1/2)\), the OAT coordinate-line set is contained in the
ball of radius \(1/2\) centered at \(b\).
The ratio of that ball's volume to the cube's volume is
\[
 r_d=\frac{\pi^{d/2}}{2^d\Gamma(d/2+1)}.
\]
Here \(\Gamma(s)=\int_0^\infty t^{s-1}e^{-t}\,dt\), \(s>0\).
In particular,
\[
 r_2=\frac{\pi}{4},\qquad r_3=\frac{\pi}{6},\qquad
 r_{12}=\frac{\pi^6}{2^{12}720}\approx0.000326,
 \qquad r_d\longrightarrow0.
\]
This is the volume of a containing region, not the volume actually
sampled by OAT. For \(d\geq2\), \(A_b\) itself has zero \(d\)-dimensional
Lebesgue measure.
\end{hbtheorem}

The centered-baseline condition is essential for this particular inscribed
ball argument. An arbitrary nominal vector need not be the cube center.
Moreover, every finite Monte Carlo design also has zero geometric volume.
Its value comes from its sampling law and integration properties;
zero sample volume by itself does not rank methods.
The interaction obstruction above is an independent information argument.

OAT remains useful for debugging a model near a known input, studying a
specified local response, or checking a justified linear or additive
representation. Its convenient common baseline also helps isolate a
particular model failure. These uses do not turn coordinate changes into
an unrestricted global analysis
\cite[Sections 2--4]{SaltelliAnnoni2010Sensitivity}.

\section{Three inexpensive alternatives}

\begin{hbdefinition}{def:p15-global-factorial}{Two-level contrasts}
Evaluate a model at every point of \(\{0,1\}^d\) after declaring the
physical low and high level of each factor. Put \(z_i=2x_i-1\) and
\(\chi_u(z)=\prod_{i\in u}z_i\), with \(\chi_\varnothing=1\).
For the corner response \(h(z)\), define
\[
 \beta_u=2^{-d}\sum_{z\in\{-1,1\}^d}h(z)\chi_u(z).
\]
The main effect of factor \(i\) is the high-level average minus the
low-level average. Under this convention it is \(2\beta_{\{i\}}\).
The usual two-factor effect is half the difference of those conditional
main effects, and is \(2\beta_{\{i,j\}}\).
All averages include the remaining corner coordinates.
\end{hbdefinition}

\begin{hbtheorem}{thm:p15-global-factorial}
 {proved=annex:p15.global.factorial}{Exact corner decomposition and its boundary}
Every corner response has the unique expansion
\[
 h(z)=\sum_{u\subseteq\{1,\ldots,d\}}\beta_u\chi_u(z).
\]
The stated main and interaction effects follow from these coefficients.
For \(d=3\), eight evaluations determine all eight coefficients.
This determines the response on the corners only: the function
\(\prod_i x_i(1-x_i)\) vanishes on all corners and is positive at the
center. A saturated corner fit has no residual degrees of freedom for
estimating an observation-noise variance without further information.
\end{hbtheorem}

\begin{hbdefinition}{def:p15-global-regression}{Standardized linear sensitivity}
Suppose \(Y,X_1,\ldots,X_d\) are square integrable, with
\(V=\operatorname{Var}(Y)>0\) and
\(\sigma_i^2=\operatorname{Var}(X_i)>0\).
Set
\[
 W=\frac{Y-EY}{\sqrt V},\qquad
 Z_i=\frac{X_i-EX_i}{\sigma_i}.
\]
Write \(Z=(Z_1,\ldots,Z_d)^\mathsf T\).
A population least-squares linear approximation is
\(\widehat W=b^\mathsf T Z\), with \(b\in\mathbb R^d\).
Let \(\Sigma=E[ZZ^\mathsf T]\) and \(c=E[ZW]\).
The explained fraction is
\(R^2=1-E[(W-\widehat W)^2]\) at the minimizing coefficient vector.
\end{hbdefinition}

\begin{hbtheorem}{thm:p15-global-regression}
 {proved=annex:p15.global.regression}{What squared regression coefficients measure}
If \(\Sigma\) is positive definite, the unique minimizing coefficient
vector is \(b=\Sigma^{-1}c\), and
\[
 R^2=b^\mathsf T\Sigma b=c^\mathsf T\Sigma^{-1}c.
\]
If the standardized inputs are uncorrelated, \(\Sigma=I\), hence
\(b_i=E[Z_iW]\) and \(R^2=\sum_i b_i^2\).
For an exactly linear additive model with independent inputs,
\[
 Y=a_0+\sum_i a_iX_i,\qquad
 b_i=\frac{a_i\sigma_i}{\sqrt V},\qquad R^2=1.
\]
The same algebra holds for an empirical least-squares fit when
standardization and covariance use the same empirical measure.
Population independence does not make finite sample columns exactly
orthogonal.
\end{hbtheorem}

A low \(R^2\) says that the chosen linear span leaves much output variation
unexplained under the chosen measure. Curvature and interactions are
possible causes; noise and omitted inputs are other causes when present.
It is not a universally calibrated amount of nonlinearity.
When the fit is poor, coefficient magnitudes describe that limited
linear projection and cannot stand for complete global factor importance.
With collinearity, individual coefficient rankings can fail to be uniquely
defined even when the fitted response is unique
\cite[Section 5.2]{SaltelliAnnoni2010Sensitivity}.

\begin{hbmethod}{def:p15-global-elementary-effects}{Elementary effects across changing backgrounds}
Normalize the declared independent input ranges or quantile coordinates
to a feasible design region. Write \(g(u)\) for the model in those
coordinates. For a feasible signed step \(\delta\ne0\), define
\[
 \operatorname{EE}_i(u;\delta)
       =\frac{g(u+\delta e_i)-g(u)}{\delta}.
\]
Choose a baseline and an ordering of the factors, then move each factor
once along a trajectory of \(d+1\) points. Repeat at \(r\) different
backgrounds. A radial alternative makes one change per factor from each
of \(r\) baselines. Independent random choice of trajectories can give
independent replicates across trajectories; separation of points alone
does not imply statistical independence.

Report the signed mean, mean absolute effect and, for \(r\geq2\), sample
spread:
\[
 \overline e_i=\frac1r\sum_{\ell=1}^r e_{i\ell},\qquad
 \overline{|e_i|}=\frac1r\sum_{\ell=1}^r|e_{i\ell}|,\qquad
 s_{e_i}^2=\frac1{r-1}\sum_{\ell=1}^r(e_{i\ell}-\overline e_i)^2.
\]
These are summaries of the specified design.
For a linear additive \(g\), every elementary effect of a given factor
equals its coefficient. Varying effects can arise from curvature,
interaction, or both; the spread does not distinguish these mechanisms.
The function \(g(u)=u_i^2\), sampled with steps from \(-1\) to \(0\) and
from \(0\) to \(1\), has signed effects \(-1,+1\): their average cancels
although the factor matters.

The nominal cost is \(r(d+1)\) evaluations, before reusing repeated
points. Saltelli and Annoni contrast two trajectories for \(d=12\),
cost \(26\), with a two-sided single-baseline OAT design costing \(25\).
Their discussion of four to ten trajectories concerns screening in the
settings cited, not a universal accuracy guarantee.
Step size, coordinates, feasible region, randomization and model failures
must be recorded \cite[Sections 5.3 and 6]{SaltelliAnnoni2010Sensitivity}.
\end{hbmethod}

\section{Variance attribution under independent inputs}

\begin{hbdefinition}{def:p15-global-variance-setting}{Product-law variance setting}
For this section, \(X_1,\ldots,X_d\) are independent, \(Y=f(X)\in L^2\),
and \(V=\operatorname{Var}(Y)>0\).
Let \(I=\{1,\ldots,d\}\) and
\(\mathcal G_u=\sigma(X_i:i\in u)\).
Write \(P_uY=E[Y\mid\mathcal G_u]\), with \(P_\varnothing Y=EY\).
All equalities between conditional versions are almost sure.

Define the functional components recursively:
\[
 Y_\varnothing=EY,\qquad
 Y_u=P_uY-\sum_{v\subsetneq u}Y_v
       =\sum_{v\subseteq u}(-1)^{|u|-|v|}P_vY.
\]
Each \(Y_u\) is a function of \(X_u\), up to a choice of measurable
version. Put \(V_u=E[Y_u^2]\) for nonempty \(u\).
\end{hbdefinition}

\begin{hbtheorem}{thm:p15-global-anova}
 {proved=annex:p15.global.anova}{Orthogonal functional decomposition}
In the product-law setting,
\[
 Y=\sum_{u\subseteq I}Y_u,\qquad
 E[Y_uY_v]=0\quad(u\ne v),\qquad
 \operatorname{Var}(Y)=\sum_{\varnothing\ne u\subseteq I}V_u.
\]
For every \(i\in u\ne\varnothing\),
\(E[Y_u\mid\mathcal G_{I\setminus\{i\}}]=0\).
These conditions uniquely determine the decomposition among square
integrable components measurable in their indicated coordinates.
\end{hbtheorem}

\begin{hbdefinition}{def:p15-global-indices}{First-order and total indices}
Saltelli and Annoni's first-order and total indices are
\[
 S_i=\frac{\operatorname{Var}(E[Y\mid X_i])}{V},\qquad
 S_{T_i}=\frac{E[\operatorname{Var}(Y\mid X_{-i})]}{V}.
\]
Here \(X_{-i}\) is the tuple of all coordinates other than \(i\).
The first numerator is the expected reduction of variance from learning
\(X_i\); the second is the expected remaining variance after learning all
the other coordinates. These are averages over the declared input law,
not promises about fixing an input at one selected nominal value.
The law of total variance in
Theorem~\ref{thm:p15-conditional-variance} gives these interpretations.
\end{hbdefinition}

\begin{hbtheorem}{thm:p15-global-indices}
 {proved=annex:p15.global.indices}{Interactions, additivity and factor fixing}
Under the product-law assumptions,
\[
 S_i=\frac{V_{\{i\}}}{V},\qquad
 S_{T_i}=\frac{\sum_{u\ni i}V_u}{V},\qquad
 0\leq S_i\leq S_{T_i}\leq1.
\]
Also
\[
 \sum_i S_i\leq1,\qquad
 1\leq\sum_iS_{T_i}\leq d.
\]
The first sum equals \(1\) exactly when the response is additive almost
surely. Equivalently, every interaction component of order at least two
vanishes, or \(S_i=S_{T_i}\) for every factor.
The total indices generally overlap; their sum is not a partition of
variance. Finally, \(S_{T_i}=0\) exactly when \(Y\) is a function of
\(X_{-i}\) almost surely under this law.
\end{hbtheorem}

Zero total effect is the appropriate variance criterion for discarding a
factor in this setting. A zero first-order effect can conceal a pure
interaction. Almost-sure factor fixing does not establish independence
from that input at every point of an enlarged domain or under a different
input law.

\begin{hbtheorem}{thm:p15-global-product-case}
 {proved=annex:p15.global.product-case}{A pure interaction missed at the origin}
Let \(X_1,X_2\) be independent and uniform on \([-1,1]\), and
\(Y=X_1X_2\). Every fixed-baseline OAT evaluation about \((0,0)\)
returns zero, but
\[
 EY=0,\qquad V=\frac19,\qquad
 S_1=S_2=0,\qquad S_{T_1}=S_{T_2}=1.
\]
The only nonzero variance component is the two-factor interaction.
\end{hbtheorem}

\begin{hbtheorem}{thm:p15-global-pick-freeze}
 {proved=annex:p15.global.pick-freeze}{Paired evaluations for the two indices}
Let \(A,B\) be independent copies of the product-law input vector, and
\(C_i=(A_i,B_{-i})\): coordinate \(i\) comes from \(A\), and each other
coordinate retains its position from \(B\). Then
\[
 \operatorname{Var}(E[Y\mid X_i])
   =E[f(A)f(C_i)]-(EY)^2,
\qquad
 E[\operatorname{Var}(Y\mid X_{-i})]
   =\frac12E[(f(B)-f(C_i))^2].
\]
With independent pairs \((A^{(\ell)},B^{(\ell)})\) across
\(\ell\), sample means of these quantities and a consistent positive
output-variance estimate give consistent ratio estimates.
Finite-sample ratios are not asserted unbiased.
For all \(d\) factors together, this construction uses \(n(d+2)\)
model evaluations, before any reuse of duplicate points.
\end{hbtheorem}

These identities are direct computational consequences of conditional
independence. For an explicit estimate, let
\(\overline y_A=n^{-1}\sum_\ell f(A^{(\ell)})\) and
\(\widehat V=n^{-1}\sum_\ell f(A^{(\ell)})^2-\overline y_A^2\).
When \(\widehat V>0\), use
\[
 \widehat S_i=
 \frac{n^{-1}\sum_\ell f(A^{(\ell)})f(C_i^{(\ell)})
                         -\overline y_A^2}{\widehat V},
\qquad
 \widehat S_{T_i}=
 \frac{(2n)^{-1}\sum_\ell
       (f(B^{(\ell)})-f(C_i^{(\ell)}))^2}{\widehat V}.
\]
Finite estimates can lie outside population bounds because of estimation
error; silently clipping them is not an uncertainty assessment.
The identities do not specify a confidence interval.
Monte Carlo error estimates require their own assumptions; deterministic
quasi-Monte Carlo error bounds require a suitable function class.
An emulator adds approximation uncertainty that must be assessed
separately. Screening first and then applying more expensive methods to
the retained factors is the source's suggested budget strategy
\cite[Sections 5.4 and 7]{SaltelliAnnoni2010Sensitivity}.

\section{The source's examples and what they establish}

The two toy functions in the source are
\[
 G(x;a)=\prod_{i=1}^d\frac{|4x_i-2|+a_i}{1+a_i},
 \qquad a_i\geq0,
\]
and
\[
\begin{split}
 D(x_1,x_2)={}&0.2e^{x_1-3}+2.2|x_2|+1.3x_2^6
              -2x_2^2-0.5x_2^4\\
 &-0.5x_1^4+2.5x_1^2+0.7x_1^3\\
 &+\frac3{(8x_1-2)^2+(5x_2-3)^2+1}
       +\sin(5x_1)\cos(3x_1^2).
\end{split}
\]
For its reported experiments the source assigns independent uniforms on
\([-1,1]\), not the customary \([0,1]\) version of the \(G\) benchmark.
It uses \(a=(0,0.01)\) for \(d=2\);
\(a=(0,0.1,0.2,0.3,0.4)\) for \(d=5\); and
\(a=(0,0.1,0.2,0.3,0.4,0.8,1)\) for \(d=7\).

\begin{hbtheorem}{thm:p15-global-source-toys}
 {proved=annex:p15.global.source-toys}{Exact calculations for the declared toy domains}
For the \(G\) function with the source's independent uniform
\([-1,1]\) inputs, set
\[
 m_i=\frac{a_i+5/2}{1+a_i},\qquad
 v_i=\frac{37}{12(1+a_i)^2}.
\]
Then
\[
 EG=\prod_i m_i,\qquad
 \operatorname{Var}(G)=
   \prod_i(m_i^2+v_i)-\prod_i m_i^2,
\qquad
 V_u=\prod_{i\in u}v_i\prod_{j\notin u}m_j^2.
\]
Consequently the first and total numerators for \(i\) are respectively
\[
 v_i\prod_{j\ne i}m_j^2,\qquad
 v_i\prod_{j\ne i}(m_j^2+v_j).
\]
For the \(D\) function,
\[
 D(1/4,3/5)-D(1/4,0)-D(0,3/5)+D(0,0)=\frac{81}{35}.
\]
Thus the mixed change contains a substantial interaction that cannot
be recovered by adding the two baseline changes.
\end{hbtheorem}

The reported experiment compares empirical output CDFs.
For the toy examples, the source describes a zero baseline with four
coordinate moves and \(1+4d\) OAT evaluations.
Its EE budget uses \(r\) trajectories with
\(r\) chosen to make \(r(d+1)\) close to \(1+4d\).
The reported counts are \(9\) versus \(9\) at \(d=2\),
\(21\) versus \(24\) at \(d=5\), and \(29\) versus \(32\) at \(d=7\).
A separate \(1000d\)-evaluation Sobol-sequence run supplies the reference
CDF. It is a numerical reference, not an exact population distribution.

The plots show the single-baseline OAT responses concentrated near the
baseline and missing much of the displayed output range, while the
illustrated EE trajectories reach a wider range. These are the reported
finite examples. They do not prove a universal EE-versus-OAT accuracy
ordering, and the empirical CDF of an arbitrary designed set is not
automatically an iid sample estimate under \(P\).

\subsection{The rate-chain comparison}
\label{sec:p15-global-rate-chain}

The application in Section 6.2 is retained as evidence for that
comparison. It uses the linear chain
\[
 \dot N_1=-\lambda_1N_1,\qquad
 \dot N_i=\lambda_{i-1}N_{i-1}-\lambda_iN_i,\quad 2\leq i\leq d,
\]
with \(N_1(0)=100\), other initial concentrations zero, and output
\(N_d(0.1\,\mathrm{s})\). Independent uncertain rates have uniform
ranges whose endpoints were selected within
\([1,100]\,\mathrm{s}^{-1}\).
The source reports the following evaluation budgets:
\[
 \begin{array}{c|rrrrrr}
 d&2&4&6&8&10&12\\ \hline
 \mathrm{OAT}&9&17&25&33&41&49\\
 \mathrm{EE}&9&15&28&36&44&52
 \end{array}
\]
It again compares with \(1000d\) quasi-Monte Carlo evaluations.
This report concerns design behavior in the given rate-chain experiment;
it does not establish a new application-domain model validation
\cite[Section 6]{SaltelliAnnoni2010Sensitivity}.

\section{Dependence, intervention and changing assumptions}

\begin{hbtheorem}{thm:p15-global-dependent-boundary}
 {proved=annex:p15.global.dependent-boundary}{Dependent inputs change the interpretation}
The two conditional-variance ratios remain defined for any joint law
with \(0<V<\infty\), and each lies in \([0,1]\).
Their independent-input ordering and decomposition need not survive.
For \(X_1=X_2=U\), where \(U\) is a fair binary variable, and \(Y=X_1\),
\[
 S_1=S_2=1,\qquad S_{T_1}=S_{T_2}=0.
\]
Both inputs individually reveal the output; either is redundant once
the other is known. Independent marginal resampling would leave the
diagonal support and analyze a different joint law.
\end{hbtheorem}

\begin{hbtheorem}{thm:p15-global-causal-boundary}
 {proved=annex:p15.global.causal-boundary}{The indices do not identify an intervention effect}
There are two structural models with identical joint distributions of
\((X_1,X_2,Y)\), and hence identical indices above, but different
distributions of \(Y\) after setting \(X_1=1\).
Thus these observational sensitivity quantities do not by themselves
identify a causal intervention response.
\end{hbtheorem}

The intervention changes a structural assignment while leaving the other
mechanisms in place; this is the distinction used in
Chapter~\ref{sec:p16-finite-causal-models}.
In deterministic simulation, changing an input of a specified causal
model can represent an intervention if that interpretation is justified.
The numerical index alone supplies no such justification.

For stochastic outputs, first specify whether the target is the mean
response \(E[Y\mid X]\) or a random realization \(Y\).
The decomposition
\[
 \operatorname{Var}(Y)
 =\operatorname{Var}(E[Y\mid X])
   +E[\operatorname{Var}(Y\mid X)]
\]
keeps input-explained variability separate from the remaining variation.
An independent random seed can instead be included as an additional
input for a decomposition of the entire stochastic output.
When \(V=0\), normalized variance indices are undefined; the conclusion
is constant output almost surely under the declared law, not universal
irrelevance on other domains.

\begin{hbtheorem}{thm:p15-global-ranking}
 {proved=annex:p15.global.ranking}{A ranking and threshold can reverse under a new range}
For \(Y=X_1+2X_2\), with independent centered inputs of variances
\(v_1,v_2>0\),
\[
 S_1=S_{T_1}=\frac{v_1}{v_1+4v_2},\qquad
 S_2=S_{T_2}=\frac{4v_2}{v_1+4v_2}.
\]
If both inputs are uniform on \([-1,1]\), these values are
\(1/5,4/5\), and \(Y\leq3\) throughout the input rectangle.
If only the range of \(X_1\) becomes \([-3,3]\), the values are
\(9/13,4/13\), the ranking reverses, and the bound \(Y\leq3\) fails,
with \(P(Y>3)=1/12\).
The derivative coefficients \(1,2\) have not changed.
More generally, over independently varying allowed variances
\(v_1\in[\ell_1,u_1]\), \(v_2\in[\ell_2,u_2]\), with positive output
variance, the inequality \(S_1\geq S_2\) holds throughout the family
exactly when \(\ell_1\geq4u_2\).
\end{hbtheorem}

\begin{hbmethod}{def:p15-global-lookup}{Choosing and reporting a sensitivity analysis}
Specify the model output and the inference that matters.
Declare feasible ranges, the joint input law, dependencies and any
alternative model formulations. Use local derivatives for a local
question, corner contrasts for the specified two-level response,
elementary effects for screening across varied backgrounds, and
variance-based indices when their assumptions and computational budget fit.
Report first and total indices together when interactions matter.

Recompute a conclusion under justified alternative ranges, dependence
structures and model assumptions. State which change reverses a ranking,
crosses a decision threshold or leaves the conclusion intact.
For sampled calculations, report design, number of model calls and
estimation uncertainty separately from mathematical identifiability.
A largest observed output is not a certified global maximum without
an additional bound or optimization argument.
A failed run is information about the model/domain combination and must
not be silently omitted. The final conclusion remains conditional on
the specified model and uncertainty description.
\end{hbmethod}
